\begin{lemma}[Lemma~2.9]
Let $\mathbb G$ be a sub-elliptic setting (\noderef{CarnotSetting}) with group law $*$ and
frame $\sigma$, let $A:\mathbb R^{m}\to\mathbb R^{m\times m}$ and
$H:\mathbb R^{m}\to\mathbb R$, let $\Om\subseteq\mathbb R^{n}$, let
$u:\mathbb R^{n}\to\mathbb R$ and let $h\in\mathbb R^{n}$. Write $L_h(x)=h*x$ and
$u^{h}:=u\circ L_h$. Then:
\begin{enumerate}
\item if $u$ is a viscosity sub-solution of $\mathcal L\cdot=0$ in $\Om$
  (\noderef{IsViscSubsol}), then $u^{h}$ is a viscosity sub-solution of
  $\mathcal L\cdot=0$ in $L_h^{-1}(\Om)=h^{-1}*\Om$;
\item if $u$ is a viscosity super-solution of $\mathcal L\cdot=0$ in $\Om$
  (\noderef{IsViscSupersol}), then $u^{h}$ is a viscosity super-solution of
  $\mathcal L\cdot=0$ in $L_h^{-1}(\Om)$.
\end{enumerate}
\end{lemma}
\begin{proof}
This is Lemma~2.9. Write $\Om':=L_h^{-1}(\Om)=\{x:\ h*x\in\Om\}$ and
$\bar u:=u^{h}=u\circ L_h$. We prove (1); (2) is the same argument with the two
inequalities reversed, and is recorded separately because \noderef{IsViscSupersol} is a
separate predicate.

\medskip\noindent\textbf{Step 0: the two cancellations and the smoothness of $L_{\pm h}$.}
For every $x$,
\[
 (-h)*(h*x)=\big((-h)*h\big)*x=0*x=x,\qquad
 h*\big((-h)*x\big)=\big(h*(-h)\big)*x=0*x=x, \tag{1}
\]
by associativity of $*$ and the identities $(-h)*h=0$, $h*(-h)=0$, $0*x=x$ of
\noderef{CarnotSetting}. In particular $L_h$ and $L_{-h}$ are mutually inverse bijections of
$\mathbb R^{n}$, and both are $C^{\infty}$, being restrictions of the $C^{\infty}$ group law
to a fixed first argument. Moreover
\[
 L_{-h}(\Om)\subseteq\Om'\qquad\text{and}\qquad L_h(\Om')\subseteq\Om , \tag{2}
\]
the first because $h*\big((-h)*z\big)=z\in\Om$ for $z\in\Om$ by (1), the second by the
definition of $\Om'$.

\medskip\noindent\textbf{Step 1: upper semi-continuity of $\bar u$ on $\Om'$.} Let
$x\in\Om'$ and let $y>\bar u(x)=u(h*x)$. Since $h*x\in\Om$ and $u$ is upper semi-continuous
on $\Om$ (first clause of \noderef{IsViscSubsol}), $u<y$ on a neighbourhood of $h*x$ in
$\Om$. The map $L_h$ is continuous and sends $\Om'$ into $\Om$ by (2), so it sends a
neighbourhood of $x$ in $\Om'$ into any prescribed neighbourhood of $h*x$ in $\Om$; hence
$\bar u=u\circ L_h<y$ on a neighbourhood of $x$ in $\Om'$. Thus $\bar u$ is upper
semi-continuous on $\Om'$.

\medskip\noindent\textbf{Step 2: the transported test function.} Let $\bar x\in\Om'$, let
$\bar\varphi:\mathbb R^{n}\to\mathbb R$ be $C^{2}$ on $\Om'$ with
$\bar x\in\argmax_{\Om'}(\bar u-\bar\varphi)$, and let $(D,B)$ be a second-order jet of
$\bar\varphi$ at $\bar x$ (\noderef{TwiceDiffAt}) with
$\mathrm{hgrad}\ \sigma\ \bar x\ D\ne0$ (\noderef{hgrad}). Put
\[
 \varphi:=\bar\varphi\circ L_{-h},\qquad\text{i.e.}\qquad \varphi(z)=\bar\varphi\big((-h)*z\big),
\]
which is the paper's $\bar\varphi\circ L_h^{-1}$ by (1). We check the three conditions that
make $\varphi$ admissible for $u$ at the point $h*\bar x\in\Om$.

\emph{(2a) $\varphi$ is $C^{2}$ on $\Om$.} $L_{-h}$ is $C^{\infty}$, hence $C^{2}$, and maps
$\Om$ into $\Om'$ by (2); $\bar\varphi$ is $C^{2}$ on $\Om'$. The composition of a $C^{2}$
map into $\Om'$ with a function $C^{2}$ on $\Om'$ is $C^{2}$ on $\Om$.

\emph{(2b) $h*\bar x\in\argmax_\Om(u-\varphi)$.} First, by (1),
$\varphi(h*\bar x)=\bar\varphi\big((-h)*(h*\bar x)\big)=\bar\varphi(\bar x)$, so
$(u-\varphi)(h*\bar x)=u(h*\bar x)-\bar\varphi(\bar x)=(\bar u-\bar\varphi)(\bar x)$. Now
let $z\in\Om$ and put $w:=(-h)*z$, which lies in $\Om'$ by (2). Maximality of
$\bar u-\bar\varphi$ over $\Om'$ at $\bar x$ gives
\[
 u(h*w)-\bar\varphi(w)=(\bar u-\bar\varphi)(w)\ \le\ (\bar u-\bar\varphi)(\bar x),
\]
and $h*w=h*\big((-h)*z\big)=z$ by (1) while $\bar\varphi(w)=\varphi(z)$ by the definition of
$\varphi$. So $(u-\varphi)(z)\le(\bar u-\bar\varphi)(\bar x)=(u-\varphi)(h*\bar x)$. As
$z\in\Om$ was arbitrary, $h*\bar x$ maximizes $u-\varphi$ over $\Om$. This is the paper's
"$u-\varphi=(u^{h}-\bar\varphi)\circ L_h^{-1}$ attains its maximum at $h*\bar y$".

\emph{(2c) the jet of $\varphi$ at $h*\bar x$, with the same horizontal data.} By (1),
$(-h)*(h*\bar x)=\bar x$, so $(D,B)$ is a second-order jet of $\bar\varphi$ at
$(-h)*(h*\bar x)$. Apply the left-invariance clause $\mathtt{leftJet}$ of
\noderef{CarnotSetting} with the group element $-h$, the base point $h*\bar x$, the function
$\bar\varphi$ and the jet $(D,B)$: it provides a second-order jet $(D',B')$ of
$y\mapsto\bar\varphi\big((-h)*y\big)=\varphi(y)$ at $h*\bar x$ with
\[
 \mathrm{hgrad}\ \sigma\ (h*\bar x)\ D'=\mathrm{hgrad}\ \sigma\ \bar x\ D,\qquad
 \mathrm{hhessS}\ \sigma\ (h*\bar x)\ D'\ B'=\mathrm{hhessS}\ \sigma\ \bar x\ D\ B
 \tag{3}
\]
(\noderef{hgrad}, \noderef{hhessS}), after rewriting $(-h)*(h*\bar x)$ as $\bar x$. This is
exactly the paper's "$\mathcal X\varphi(h*\bar y)=\mathcal X\bar\varphi(\bar y)$", together
with the corresponding identity for the symmetrized horizontal second derivative; it is the
content of the left-invariance of the fields $X_i$, i.e. of
$X_i(\varphi\circ L_h)=(X_i\varphi)\circ L_h$, read on $2$-jets as pm-0002 prescribes.

In particular $\mathrm{hgrad}\ \sigma\ (h*\bar x)\ D'=\mathrm{hgrad}\ \sigma\ \bar x\ D\ne0$.

\medskip\noindent\textbf{Step 3: conclusion.} By (2a), (2b), (2c) the function $\varphi$ is
admissible for \noderef{IsViscSubsol} at the point $h*\bar x\in\Om$ with the jet $(D',B')$,
so
\[
 \mathrm{Lop}\ \sigma\ A\ H\ (h*\bar x)\ D'\ B'\ \le\ 0\qquad(\noderef{Lop}).
\]
By \noderef{Lop}, $\mathrm{Lop}\ \sigma\ A\ H\ x\ D\ B$ depends on $(x,D,B)$ only through
$\mathrm{hgrad}\ \sigma\ x\ D$ and $\mathrm{hhessS}\ \sigma\ x\ D\ B$, namely
$\mathrm{Lop}=-\Tr\big(A(\mathrm{hgrad})\,\mathrm{hhessS}\big)+H(\mathrm{hgrad})$. Hence
(3) gives
\[
 \mathrm{Lop}\ \sigma\ A\ H\ (h*\bar x)\ D'\ B'=\mathrm{Lop}\ \sigma\ A\ H\ \bar x\ D\ B ,
\]
and therefore $\mathrm{Lop}\ \sigma\ A\ H\ \bar x\ D\ B\le0$. Since $\bar x$, $\bar\varphi$
and $(D,B)$ were arbitrary, and $\bar u$ is upper semi-continuous on $\Om'$ by Step 1,
$\bar u=u^{h}$ is a viscosity sub-solution of $\mathcal L\cdot=0$ in $\Om'$ in the sense of
\noderef{IsViscSubsol}. This proves (1).

\medskip\noindent\textbf{Step 4: the super-solution half.} For (2), Step 0 and Steps (2a),
(2c) do not involve $u$ and are verbatim. Step 1 becomes lower semi-continuity of $\bar u$
on $\Om'$, by the same pull-back of a neighbourhood through the continuous $L_h$. In (2b)
the inequality reverses: with $w=(-h)*z\in\Om'$, minimality of $\bar u-\bar\varphi$ over
$\Om'$ at $\bar x$ gives $(\bar u-\bar\varphi)(\bar x)\le(\bar u-\bar\varphi)(w)$, i.e.
$(u-\varphi)(h*\bar x)\le(u-\varphi)(z)$, so $h*\bar x\in\argmin_\Om(u-\varphi)$. Feeding
$\varphi$ and $(D',B')$ to \noderef{IsViscSupersol} gives
$0\le\mathrm{Lop}\ \sigma\ A\ H\ (h*\bar x)\ D'\ B'=\mathrm{Lop}\ \sigma\ A\ H\ \bar x\ D\ B$,
which is (2).
\end{proof}
